\MathBookSourcePage{341}
\subsection{抽象混合形式}
\label{sec:ch12-abstract-mixed-formulation}
\setcounter{thmcount}{0}
\setcounter{equation}{0}

现在抽取上述两个问题的关键特征. 设 $V$ 和 $\Pi$ 是两个 Hilbert 空间, 且有两个双线性形式 $a(\cdot,\cdot):V\times V\to\mathbb R$ 和 $b(\cdot,\cdot):V\times\Pi\to\mathbb R$. 自然假设这些形式连续:
\MBSourceItem{1}
\begin{equation}
\MathBookFormulaID{formula-ch12-bilinear-form-continuity}
\label{eq:ch12-bilinear-form-continuity}
|a(u,v)|\leq C\|u\|_V\|v\|_V,
\qquad
|b(v,p)|\leq C\|v\|_V\|p\|_\Pi.
\end{equation}
暂时推迟讨论适当的强制性形式.

$V$ 与 $\Pi$ 之间没有唯一关系, 但有一个算子提供联系. 假设 $D:V\to\Pi$ 连续,
\MBSourceItem{2}
\begin{equation}
\MathBookFormulaID{formula-ch12-D-continuity}
\label{eq:ch12-D-continuity}
\|Dv\|_\Pi\leq C\|v\|_V,
\end{equation}
并作简化假设
\MBSourceItem{3}
\begin{equation}
\MathBookFormulaID{formula-ch12-b-as-D-inner-product}
\label{eq:ch12-b-as-D-inner-product}
b(v,p)=(Dv,p)_\Pi.
\end{equation}

\MathBookSourcePage{342}
考虑如下变分问题: 求 $u\in V$ 和 $p\in\Pi$, 使得
\MBSourceItem{4}
\begin{equation}
\MathBookFormulaID{formula-ch12-general-mixed-problem}
\label{eq:ch12-general-mixed-problem}
\begin{aligned}
a(u,v)+b(v,p)&=F(v)&&\forall v\in V,\\
b(u,q)&=G(q)&&\forall q\in\Pi,
\end{aligned}
\end{equation}
其中 $F\in V'$ 且 $G\in\Pi'$. 还存在一些替代形式(Ciarlet 与 Lions 1991), 其中存在算子 $D':\Pi\to V$ 且 $b(v,p)=(v,D'p)_V$.

通过
\MBSourceItem{5}
\begin{equation}
\MathBookFormulaID{formula-ch12-kernel-Z}
\label{eq:ch12-kernel-Z}
Z:=\{v\in V:b(v,q)=0\quad\forall q\in\Pi\}
\end{equation}
定义 $V$ 的闭子空间. 暂设~\eqref{eq:ch12-general-mixed-problem} 中 $G=0$. 则~\eqref{eq:ch12-general-mixed-problem} 确定 $u$ 等价于求 $u\in Z$, 使得
\MBSourceItem{6}
\begin{equation}
\MathBookFormulaID{formula-ch12-reduced-kernel-problem}
\label{eq:ch12-reduced-kernel-problem}
a(u,v)=F(v)\qquad\forall v\in Z.
\end{equation}
只要 $a(\cdot,\cdot)$ 强制, 即
\MBSourceItem{7}
\begin{equation}
\MathBookFormulaID{formula-ch12-coercivity-on-Z}
\label{eq:ch12-coercivity-on-Z}
\alpha\|v\|_V^2\leq a(v,v)\qquad\forall v\in Z,
\end{equation}
由 Lax--Milgram 定理~\MBUnresolvedReference{thm:ch02-source-2-7-7}{2.7.7}, 该问题就是适定的.

强制性条件~\eqref{eq:ch12-coercivity-on-Z} 对上一节两个问题都成立. 对 Stokes 形式~\eqref{eq:ch12-stokes-a-form}, 它对所有 $v\in V$ 成立; 但对其他形式(参见习题~\ref{exr:ch12-korn-variant-coercivity}), 它只在 $Z$ 上成立. 对标量椭圆问题, 必须允许限制到 $Z$, 在该情形
\[
\MathBookFormulaID{formula-ch12-hdiv-divergence-free-kernel}
Z=\{\chvec v\in H(\operatorname{div}):\operatorname{div}\chvec v=0\}.
\]
对 $\chvec v\in Z$, $\|\chvec v\|_{H(\operatorname{div})}=\|\chvec v\|_{L^2(\Omega)^n}$, 强制性由系数 $a_{ij}$ 有界得到.

若~\eqref{eq:ch12-general-mixed-problem} 中 $G\neq0$, 可约化到 $G=0$ 的情形. 设 $u_0\in V$ 是 $b(u_0,q)=G(q)$ 的任一解. 对上一节的例子, 这相当于求解 $\operatorname{div}u_0=g$, 其可通过多种方式完成(参见引理~\ref{lem:ch11-divergence-right-inverse} 以及 Arnold, Scott 与 Vogelius 1988). 此时~\eqref{eq:ch12-general-mixed-problem} 的解形为 $u=u_1+u_0$, 其中 $u_1\in Z$ 满足
\[
\MathBookFormulaID{formula-ch12-shifted-kernel-problem}
a(u_1,v)=F(v)-a(u_0,v)\qquad\forall v\in Z.
\]
这仍具有~\eqref{eq:ch12-reduced-kernel-problem} 的形式. 因而本节其余部分集中于
\MBSourceItem{8}
\begin{equation}
\MathBookFormulaID{formula-ch12-homogeneous-mixed-problem}
\label{eq:ch12-homogeneous-mixed-problem}
\begin{aligned}
a(u,v)+b(v,p)&=F(v)&&\forall v\in V,\\
b(u,q)&=0&&\forall q\in\Pi,
\end{aligned}
\end{equation}
一般情形将在第~\ref{sec:ch12-discrete-inf-sup} 节回到.

\MathBookSourcePage{343}
由~\eqref{eq:ch12-reduced-kernel-problem} 适定地确定 $u$ 后, 再求 $p\in\Pi$, 使得
\MBSourceItem{9}
\begin{equation}
\MathBookFormulaID{formula-ch12-pressure-recovery}
\label{eq:ch12-pressure-recovery}
b(v,p)=-a(u,v)+F(v)\qquad\forall v\in V.
\end{equation}
该问题的适定性来自一种新的强制性. 为说明这一新条件, 重新考察~\eqref{eq:ch12-coercivity-on-Z}. 该条件蕴含
\[
\MathBookFormulaID{formula-ch12-coercivity-sup-form}
\alpha\|v\|_V\leq\sup_{w\in V}\frac{a(v,w)}{\|w\|_V}
\qquad\forall v\in V
\]
(取 $w=v$), 后一条件足以给出涉及 $a(\cdot,\cdot)$ 的估计. 类似地, 对 $b(\cdot,\cdot)$ 考虑条件
\MBSourceItem{10}
\begin{equation}
\MathBookFormulaID{formula-ch12-continuous-inf-sup}
\label{eq:ch12-continuous-inf-sup}
\beta\|p\|_\Pi\leq
\sup_{v\in V}\frac{b(v,p)}{\|v\|_V}
\qquad\forall p\in\Pi.
\end{equation}
问题~\eqref{eq:ch12-pressure-recovery} 具有形式
\MBSourceItem{11}
\begin{equation}
\MathBookFormulaID{formula-ch12-abstract-b-equation}
\label{eq:ch12-abstract-b-equation}
b(v,p)=\widetilde F(v)\qquad\forall v\in V,
\end{equation}
其中 $\widetilde F(v)=0$ 对所有 $v\in Z$ 成立. 条件~\eqref{eq:ch12-continuous-inf-sup} 蕴含解的唯一性. 解的存在性还需进一步说明.

\MBSourceItem{12}
\begin{Lemma}
\label{lem:ch12-inf-sup-unique-solvability}
假设~\eqref{eq:ch12-bilinear-form-continuity} 和~\eqref{eq:ch12-continuous-inf-sup} 对某个 $\beta>0$ 成立. 则~\eqref{eq:ch12-abstract-b-equation} 有唯一解.
\end{Lemma}

\begin{Proof}
令 $Z^\perp$ 表示 $V$ 中 $Z$ 的正交补. 由于 $b(\cdot,\cdot)$ 在 $Z$ 上的行为平凡, 可把注意力限制到 $v\in Z^\perp$. 给定 $p\in\Pi$, 线性形式 $v\mapsto b(v,p)$ 在 $Z^\perp$ 上连续, 故 Riesz 表示定理~\MBUnresolvedReference{thm:ch02-source-2-4-2}{2.4.2} 保证存在 $Tp\in Z^\perp$, 使得
\MBSourceItem{13}
\begin{equation}
\MathBookFormulaID{formula-ch12-operator-T-definition}
\label{eq:ch12-operator-T-definition}
(Tp,v)_V=b(v,p)\qquad\forall v\in Z^\perp.
\end{equation}
此外, $T$ 线性, 且由~\eqref{eq:ch12-bilinear-form-continuity}
\[
\MathBookFormulaID{formula-ch12-T-continuity}
\|Tp\|_V=\sup_{v\in Z^\perp}\frac{b(v,p)}{\|v\|_V}
\leq C\|p\|_\Pi.
\]
令 $R$ 表示 $T$ 在 $Z^\perp$ 中的像. 若能证明 $R=Z^\perp$, 则再次应用 Riesz 表示定理即可完成证明, 因为 $\widetilde F$ 总可表示为 $\widetilde F(v)=(u,v)_V$.

\MathBookSourcePage{344}
对某个 $u\in Z^\perp$, 只需取满足 $Tp=u$ 的 $p$. 为证明 $R=Z^\perp$, 先证明 $R$ 闭. 设 $\{p_j\}\subset\Pi$ 且 $Tp_j\to w\in Z^\perp$. 则
\[
\MathBookFormulaID{formula-ch12-T-range-closed-estimate}
\beta\|p_j-p_k\|_\Pi
\leq\sup_{v\in Z^\perp}\frac{b(v,p_j-p_k)}{\|v\|_V}
=\sup_{v\in Z^\perp}\frac{(v,Tp_j-Tp_k)_V}{\|v\|_V}
=\|Tp_j-Tp_k\|_V.
\]
故 $\{p_j\}$ 在 $\Pi$ 中为 Cauchy 序列. 令 $q=\lim p_j$. 由 $T$ 的连续性, $Tq=w$, 因而 $R$ 闭. 若 $R\neq Z^\perp$, 取 $0\neq v\in R^\perp$. 则 $b(v,q)=(v,Tq)_V=0$ 对所有 $q\in\Pi$ 成立. 这意味着 $v\in Z$, 与 $v\in Z^\perp$ 矛盾. 所以 $R=Z^\perp$.
\end{Proof}

条件~\eqref{eq:ch12-continuous-inf-sup} 对上一节引入的两个问题都成立, 且证明相似. 由引理~\ref{lem:ch11-divergence-right-inverse}, 可解欠定系统 $\operatorname{div}\chvec v=p$, 并取 $\chvec v\in H^1(\Omega)^n$ 满足 $\|\chvec v\|_{H^1(\Omega)^n}\leq(1/\beta)\|p\|_{L^2(\Omega)}$. 若进一步 $\int_\Omega p\,dx=0$, 可取 $\chvec v\in\mathring H^1(\Omega)^n$. 因此对上一节研究的两种情形之一, 都存在 $\chvec v\in V$, 使得
\[
\MathBookFormulaID{formula-ch12-inf-sup-divergence-verification}
\beta\|p\|_{L^2(\Omega)}
=\beta\frac{b(\chvec v,p)}{\|p\|_{L^2(\Omega)}}
\leq\frac{b(\chvec v,p)}{\|\chvec v\|_{H^1(\Omega)^n}}
\leq\sqrt n\,\frac{b(\chvec v,p)}{\|\chvec v\|_{H(\operatorname{div})}},
\]
这验证了两种情形的~\eqref{eq:ch12-continuous-inf-sup}.
